\subsection{Marco Teórico y conceptual}

\subsubsection{Fundamento teórico del proyecto de Innovación / Mejora / Creatividad}

\paragraph{Sistema web}

\paragraph{}

\paragraph{}

\paragraph{}

\subsubsection{Conceptos y Términos utilizados}
\paragraph{Gestión de Nóminas}
Según Quispe Noa \cite{quispe2025implementacion}, la gestión de nómina abarca el cálculo, la organización y la ejecución de pagos por remuneraciones al empleado, incluyendo bonos de acuerdo a ley, aportaciones y descuentos.

\paragraph{Cálculo de Planillas}
Lima Trujillo \cite{lima2022implementacion} señala que el cálculo de planillas permite a la empresa llevar un registro de las remuneraciones dadas al empleado.

\paragraph{Retenciones y Aportes}
Aguilar Guizado y Sotelo Herrera \cite{aguilar2022desarrollo} mencionan que las retenciones y aportaciones son el descuento que se realiza con base en el ingreso del empleado al realizar el cálculo base para derivar el descuento al régimen de pensiones o tipo de renta al cual el empleado pertenece.

\paragraph{PostgreSQL}
Es un motor de base de datos desarrollado por la Universidad de California en Berkeley, creado para la gestión de datos estructurados usando el lenguaje SQL, utilizando como identificador único el UUID v7.

\paragraph{PrismaORM}
Es un mapeador de datos que sirve como capa intermedia entre el programa y la base de datos, funcionando como puente de abstracción en vez de escribir consultas SQL directamente.

\paragraph{Redis}
Es un motor de base de datos en memoria. En lugar de guardar los datos en discos duros o en un componente de hardware tradicional, utiliza la memoria de acceso rápido para obtener tiempos de respuesta rápidos, funcionando con una estructura de clave-valor.

\paragraph{NestJS}
Framework orientado al desarrollo de aplicaciones backend altamente escalables de código abierto, cuya arquitectura base es la arquitectura modular, lo cual permite desarrollar aplicaciones escalables y mantenibles.

\paragraph{Beneficios de empleado}
Son los beneficios que la empresa debe dar a los empleados, bajo el régimen regulatorio de la nación.

\paragraph{CTS}
Beneficio social al cual el empleado accede, siendo el principal aportador el empleador; sirve como seguro de desempleo para el empleado.

\paragraph{AFP}
Régimen de pensiones privado, donde el empleado aporta alrededor del 13\% de su salario bruto (incluyendo la comisión por mantener la rentabilidad de los aportes), gestionado en una cuenta individual de capitalización, para tener un saldo de jubilación que abarca todas las aportaciones realizadas durante la vida laboral del empleado.

\paragraph{ONP}
Régimen de pensiones público, donde el empleado aporta el 13\% de su salario bruto a una cuenta colectiva dependiente del tiempo de aportes, que es a partir de 10 a 20 años, a diferencia de la AFP, que es un fondo individual acumulado.

\paragraph{Gratificaciones}
Beneficio social obligatorio que es pagado por el empleador a sus trabajadores, dado dos veces al año en los meses de julio y diciembre; este beneficio está libre de descuentos y es equivalente a un sueldo mensual completo.

\paragraph{Asignación Familiar}
Es el beneficio social orientado a apoyar a los trabajadores del sector privado que tienen a su cargo familiares (hijos a su cargo), correspondiente al 10\% de la remuneración mínima vital.

\paragraph{API}
Conjunto de reglas, protocolos y herramientas que permite la comunicación entre sistemas informáticos.

\paragraph{API REST}
Tipo de API orientado a la comunicación web, siguiendo reglas arquitectónicas de representación de transferencia de estado.

\paragraph{ViteJs}
Herramienta de desarrollo y empaquetado para aplicativos frontend.

\paragraph{JSON}
Formato de organización para intercambio y envío de datos, cuya estructura es del tipo clave-valor, listados o arreglos.

\paragraph{Arquitectura Clean}
Metodología de diseño que separa la lógica de negocio de los detalles tecnológicos, dando independencia al proyecto respecto al tipo de tecnología que se aplique.

\paragraph{Encriptado de datos}
Proceso por el cual la información legible se transforma, utilizando un algoritmo y una clave, a un formato ilegible; en este caso, el cifrado tipo AES-256-GCM.

\paragraph{Herramientas de documentación de API}
Herramientas de documentación para aplicaciones que permiten normalizar y visualizar las aplicaciones, siendo Swagger/OpenAPI la documentación e interacción de los endpoints.

\paragraph{Periodo}
Bloques de tiempo exactos en los cuales están vigentes determinantes laborales, tales como la remuneración mínima, UIT, etc.

\paragraph{Testeo de API}
Herramientas centradas en la validación, simulación y verificación de resultados de una aplicación, tales como Postman para probar peticiones HTTP en tiempo real, y Jest, un framework de testeo para la ejecución de pruebas unitarias.

\paragraph{JWT}
Herramienta de validación de usuario en cada petición hecha por este, mediante tokens firmados sin que esa parte recaiga en la base de datos.

\paragraph{RMV}
Remuneración mínima vital, establecida por el Estado peruano, la cual es el mínimo que un empleador del sector privado debe pagar a un empleado.

\paragraph{UIT}
Valor referencial basado en soles, que es utilizado por el Estado peruano para estandarizar montos en el área legal y tributaria.

\paragraph{SUNAT}
Agente regulador del Estado en cuanto a retención, recaudación y administración en las operaciones tributarias dentro del país.

\paragraph{SUNAFIL}
Supervisor del organismo público en el cumplimiento del régimen laboral.

\paragraph{Régimen de salud público y privado}
Es el conjunto de normativas que garantizan la cobertura médica del trabajador entre los distintos sistemas de salud, como EsSalud, EPS o un seguro pagado por la empresa.

\paragraph{Derechohabientes}
Familiar directo del trabajador que obtiene el derecho de cobertura del seguro de salud del empleado y otros beneficios del empleado dados por la entidad donde trabaja.
